\subsection{相对不变测度}

命题 10. --- 设 G 为局部紧群，\(\mu\) 为 G 上乘子为 \(\chi\) 的相对左不变测度。如果 \(\chi_{1}\) 是 G 在 \(\mathbf{C}^{*}\) 中的连续表示，则测度 \(\chi_{1} \cdot \mu\) 是乘子为 \(\chi_{1}\chi\) 的相对左不变测度。

这是因为，

\[
\begin{array}{r} \pmb {\gamma} (s) (\chi_ {1} \cdot \mu) = (\pmb {\gamma} (s) \chi_ {1}) \cdot (\pmb {\gamma} (s) \mu) = (\chi_ {1} (s ^ {- 1}) \chi_ {1}) \cdot (\chi (s) ^ {- 1} \mu) \\ = (\chi_ {1} \chi) (s) ^ {- 1} (\chi_ {1} \cdot \mu). \end{array}
\]

推论 1. --- 设 \(\mu\) 是 G 上的左哈尔测度。G 上的非零测度 \(\nu\) 相对左不变，当且仅当它具有 \(a\chi \cdot \mu\) 的形式，其中 \(a \in \mathbf{C}^*\)，且 \(\chi\) 是 G 在 \(\mathbf{C}^*\) 中的连续表示；此时其乘子为 \(\chi\)。

条件是充分的（命题 10）。另一方面，如果 \(\nu\) 是乘子为 \(\chi\) 的非零相对左不变测度，则 \(\chi^{-1} \cdot \nu\) 是左不变的（命题 10），因此具有 \(a\mu\) 的形式，其中 \(a \in \mathbf{C}^*\)（第 2 号，定理 1 的推论）。

推论 2. --- 每个相对左不变测度都是相对右不变的。

这是因为，在推论 1 的记号下，

\[
\begin{array}{c} \boldsymbol {\delta} (s) (\chi \cdot \mu) = \bigl (\boldsymbol {\delta} (s) \chi \bigr) \cdot \bigl (\boldsymbol {\delta} (s) \mu \bigr) = \bigl (\chi (s) \chi \bigr) \cdot \bigl (\Delta_ {\mathrm{G}} (s) \mu \bigr) \\ = (\chi \Delta_ {\mathrm{G}}) (s) (\chi \cdot \mu). \end{array}\tag{34}
\]

由于推论 2，以下谈到 G 上的相对不变测度时不再另加说明。相对不变测度的特殊情形包括左哈尔测度和右哈尔测度。给定 G 上的相对不变测度 \(\nu\)，区分其左乘子 \(\chi\) 和右乘子 \(\chi'\) 很方便，它们定义为 \(\pmb{\gamma}(s)\nu = \chi(s)^{-1}\nu\)、\(\pmb{\delta}(s)\nu = \chi'(s)\nu\)。由（34），这些乘子满足关系

\[
\chi^ {\prime} = \chi \Delta_ {\mathrm{G}}.\tag{35}
\]

仍记左哈尔测度为 \(\mu\)，有

\[
\check {\nu} = (\chi \cdot \mu) ^ {\vee} = \check {\chi} \cdot \check {\mu} = (\chi^ {- 1} \Delta_ {\mathrm{G}} ^ {- 1}) \cdot \mu ,
\]

因此 \(\check{\nu}\) 是左乘子为 \(\chi^{-1}\Delta_{\mathrm{G}}^{-1}\) 且右乘子为 \(\chi^{-1}\) 的相对不变测度。

对于每个相对不变测度，可忽略函数、局部可忽略函数、可测函数和局部可积函数的概念都相同。

